By full row rank, the \(p\) rows of \(C\) are independent. Extend them to the rows of a matrix \(T_0\in\operatorname{GL}_n(\mathbb R)\), preserving \((T_0)_{\mathsf O,:}=C\), and put \(W_0=T_0^{-1}\). The \((x,x)\)-entry of \(T_0W_0=I_n\) is
\[
 1=\sum_{j=1}^n C_{xj}(W_0)_{jx}. \tag{10}
\]
Thus some \(j\) has both \(C_{xj}\ne0\) and \((W_0)_{jx}\ne0\). The inverse-cofactor formula says that this latter inequality is equivalent to \(\det((T_0)_{-x,-j})\ne0\). Its \(p-1\) observed rows other than \(x\), namely the rows of \(C_{-x,-j}\), are therefore independent. Hence \(j\in J_x(C)\), proving nonemptiness.

Now fix any \(j\in J_x(C)\). We first justify the shear prescribed by \(\operatorname{Cert}\). Because the rows of \(C_{-x,:}\) are independent,
\[
 \operatorname{rank}(C_{-x,-j})<p-1
 \quad\Longleftrightarrow\quad
 e_j^\top\in\operatorname{rowspan}(C_{-x,:}). \tag{11}
\]
Indeed, a nonzero dependence after deletion gives a nonzero linear combination of the independent full rows supported only at coordinate \(j\), and hence a nonzero multiple of \(e_j^\top\); the converse follows by deleting coordinate \(j\). Row \(j\) of \(W_0\) is the unique coordinate vector of \(e_j^\top\) in the row basis of \(T_0\), since
\[
 e_j^\top=(e_j^\top W_0)T_0.
\]
Consequently (11) is also equivalent to
\[
 (W_0)_{jx}=0
 \quad\text{and}\quad
 (W_0)_{j\ell}=0\ \text{for every }\ell\in\mathsf L. \tag{12}
\]
The admissibility of \(j\) rules out (12). If \((W_0)_{jx}\ne0\), set \(T_j=T_0\) and \(W_j=W_0\). Otherwise (12) supplies \(\ell\in\mathsf L\) with \((W_0)_{j\ell}\ne0\). Since \(x\ne\ell\), the matrices \(I_n+e_\ell e_x^\top\) and \(I_n-e_\ell e_x^\top\) are inverse, so the prescribed shear satisfies
\[
 T_j=(I_n+e_\ell e_x^\top)T_0,
 \qquad
 W_j=W_0(I_n-e_\ell e_x^\top),
 \qquad
 (W_j)_{jx}=-(W_0)_{j\ell}\ne0. \tag{13}
\]
Only latent row \(\ell\) changes, hence \((T_j)_{\mathsf O,:}=C\), and \(T_j\) remains invertible.

Because \(T_j=W_j^{-1}\) and \((T_j)_{xj}=C_{xj}\ne0\), another inverse-cofactor identity gives
\[
 \det((W_j)_{-j,-x})\ne0. \tag{14}
\]
The determinant expansion in (14) contains a nonzero product. Its indices give a perfect matching from the source rows other than \(j\) to the equation columns other than \(x\). Augment it by assigning row \(j\) to column \(x\), and let \(\Pi_j\) permute the rows of \(W_j\) so that every diagonal entry of \(\Pi_jW_j\) is nonzero and row \(x\) is the old row \(j\). Set
\[
 D_j=\operatorname{diag}(\Pi_jW_j),
 \qquad Q_j=D_j^{-1}\Pi_jW_j,
 \qquad H_j=D_j^{-1}\Pi_j. \tag{15}
\]
Then \(D_j\) is invertible, \(Q_j\) is invertible with unit diagonal, and \(H_j\) is monomial. Moreover,
\[
 Q_j^{-1}H_j
 =W_j^{-1}\Pi_j^{-1}D_jD_j^{-1}\Pi_j
 =T_j, \tag{16}
\]
so \([Q_j^{-1}H_j]_{\mathsf O,:}=C\) and \(\mathcal L(C\varepsilon)=P_X\). The canonical sources \(\varepsilon_1,\ldots,\varepsilon_n\) are mutually independent, nonconstant, and non-Gaussian by the population assumptions, while the columns of \(C\) are nonzero and pairwise nonproportional. Thus all irreducible LiNG and observed-law requirements in the definition of \(\mathfrak F_x(P_X)\) hold.

The augmentation in (15) gives \((H_j)_{xj}\ne0\), so \(s_{\mathcal M_j}(x)=j\). From (16),
\[
 C_{xj}=(Q_j^{-1})_{xx}(H_j)_{xj}.
\]
Both \(C_{xj}\) and \((H_j)_{xj}\) are nonzero; hence \((Q_j^{-1})_{xx}\ne0\). The cofactor formula and invertibility of \(Q_j\) now give \(\det((Q_j)_{-x,-x})\ne0\). Therefore \(\mathcal M_j=(n,Q_j,H_j,\varepsilon)\) belongs to \(\mathfrak F_x(P_X)\). Applying the law-fiber response-rank lemma to this completion and using (16) proves both displayed response formulas.

For the operation count, Gaussian elimination extends the row basis and inverts \(T_0\) in \(O(n^3)\) arithmetic operations. Equivalence (12) classifies every deletion rank by scanning \(W_0\), in \(O(n^2)\) operations after the inverse. Each of at most \(n\) witnesses uses one \(O(n^2)\) shear, at most \(n\) augmenting-path searches each scanning at most \(n^2\) potential matching edges, and \(O(n^2)\) row normalization. Hence all witnesses require \(O(n^4)\) operations. This count uses exact zero tests in the real-arithmetic model; it is not a floating-point or bit-complexity statement. If \(C\) is rational, Gaussian elimination can choose a rational row-basis extension, and inversion, the integral shear, permutation, and divisions by the nonzero rational diagonal entries in (15) preserve rationality.